% Generated from locale/tr/content/first-order-logic/natural-deduction/proof-theoretic-notions.tex; do not hand-edit.


\iftag{FOL}
      {\olfileid[tr]{fol}{ntd}{ptn}}
      {\olfileid[tr]{pl}{ntd}{ptn}}

\olsection{İspat Kuramsal Kavramlar}

\begin{editorial}
Bu bölüm, doğal tümdengelim için ispatlanabilirlik bağıntısı ile tutarlılığın
tanımlarını bir araya getirir.
\end{editorial}

\begin{explain}
Bir dizi önemli semantik kavramı (geçerlilik, mantıksal gerektirme,
sağlanabilirlik) tanımladığımız gibi şimdi bunlara karşılık gelen
\emph{ispat kuramsal kavramları} tanımlıyoruz. Bunlar, !!{sentence}s ile
!!{structure}s arasındaki sağlama bağıntısına değil, belirli
!!{derivability} veya !!{nonderivability} durumlarına; yani birtakım
!!{sentence}s ile diğerleri arasındaki türetim bağıntısına başvurularak
tanımlanır. Bu kavramların çakıştığının keşfedilmesi
önemliydi. Çakıştıkları iddiası \emph{sağlamlık} ve \emph{tamlık
teoremlerinin} içeriğidir.
\end{explain}


\begin{defn}[Teoremler]
!!{sentence}~$!A$, doğal tümdengelimde $!A$ için !!a{derivation} varsa ve bu
türetimde bütün varsayımlar !!{discharged} ise \emph{teorem}dir.
$!A$ teoremse
$\Proves !A$, değilse $\Proves/ !A$ yazarız.
\end{defn}

\begin{defn}[!!^{derivability}]
!!^a{sentence} $!A$ için \emph{!!{derivable}} terimini, kaynak bir
!!{sentence}s kümesi~$\Gamma$ olduğunda şu koşulla kullanırız: sonucu $!A$
olan bir !!{derivation} bulunmalı ve her varsayım ya !!{discharged} olmalı ya
da $\Gamma$ içinde bulunmalıdır. Bunu $\Gamma \Proves !A$ ile gösteririz.
$!A$, $\Gamma$'dan !!{derivable} değilse $\Gamma \Proves/ !A$ yazarız.
\end{defn}

\begin{defn}[Tutarlılık]
!!{sentence}s kümesi~$\Gamma$, ancak ve ancak $\Gamma \Proves \lfalse$ ise
\emph{tutarsızdır}. $\Gamma$ tutarsız değilse, yani $\Gamma \Proves/
\lfalse$ ise ona \emph{tutarlı} deriz.
\end{defn}

\begin{prop}[Yansıma]
\ollabel{prop:reflexivity}
$!A \in \Gamma$ ise $\Gamma \Proves !A$ olur.
\end{prop}

\begin{proof}
$!A$ varsayımı tek başına $!A$ için !!a{derivation} olur; buradaki her
!!{undischarged} varsayım (yani $!A$) $\Gamma$ içindedir.
\end{proof}
  

\begin{prop}[Monotonluk]
\ollabel{prop:monotonicity}
$\Gamma \subseteq \Delta$ ve $\Gamma \Proves !A$ ise $\Delta \Proves !A$
olur.
\end{prop}

\begin{proof}
$\Gamma$'dan $!A$ için her !!{derivation}, aynı zamanda $\Delta$'dan $!A$
için !!a{derivation} olur.
\end{proof}

\begin{prop}[Geçişlilik]
\ollabel{prop:transitivity}
$\Gamma \Proves !A$ ve $\{!A\} \cup \Delta \Proves !B$ ise
$\Gamma \cup \Delta \Proves !B$ olur.
\end{prop}

\begin{proof}
$\Gamma \Proves !A$ ise $!A$ için !!a{derivation}~$\delta_0$ vardır; bunun
bütün !!{undischarged} varsayımları $\Gamma$ içindedir.
$\{!A\} \cup \Delta \Proves !B$ ise $!B$ için !!a{derivation}~$\delta_1$
vardır; bunun bütün !!{undischarged} varsayımları $\{!A\} \cup \Delta$
içindedir. Şimdi şunu ele alalım:
\begin{prooftree}
  \AxiomC{$\Delta, \Discharge{!A}{1}$}
  \RightLabel{$\delta_1$}
  \DeduceC{$!B$}
  \DischargeRule{\Intro{\lif}}{1}
  \UnaryInfC{$!A \lif !B$}
  \AxiomC{$\Gamma$}
  \RightLabel{$\delta_0$}
  \DeduceC{$!A$}
  \RightLabel{\Elim{\lif}}
  \BinaryInfC{$!B$}
\end{prooftree}
Artık !!{undischarged} varsayımların hepsi $\Gamma \cup \Delta$ içindedir;
dolayısıyla bu, $\Gamma \cup \Delta \Proves !B$ olduğunu gösterir.
\end{proof}

$\Gamma = \{!A_1, !A_2, \ldots, !A_k\}$ sonlu bir küme olduğunda
$\Gamma \Proves !B$ yerine sadeleştirilmiş $!A_1,!A_2,\ldots,!A_k \Proves
!B$ gösterimini kullanabiliriz. Özellikle $!A \Proves !B$,
$\{!A\} \Proves !B$ demektir.

$\Gamma \Proves !A$ ve $!A \Proves !B$ ise $\Gamma \Proves !B$ olduğuna
dikkat edin. Ayrıca $!A_1, \dots, !A_n \Proves !B$ ve her~$i$ için
$\Gamma \Proves !A_i$ ise $\Gamma \Proves !B$ olduğu da çıkar.

\begin{prop}
\ollabel{prop:incons}
Aşağıdakiler eşdeğerdir.
    \begin{enumerate}
      \item \( \Gamma \) tutarsızdır.
      \item Her !!{sentence}~\( {!A} \) için \( \Gamma \Proves {!A} \).
      \item Bazı !!{sentence}~\( {!A} \) için hem \( \Gamma \Proves {!A} \)
      hem de \( \Gamma \Proves \lnot {!A} \).
    \end{enumerate}
\end{prop}

\begin{proof}
Alıştırma.
\end{proof}

\tagprob{FOL}
\begin{prob}
\olref[fol][ntd][ptn]{prop:incons} önermesini ispatlayın.
\end{prob}
\tagendprob

\tagprob{notFOL}
\begin{prob}
\olref[pl][ntd][ptn]{prop:incons} önermesini ispatlayın.
\end{prob}
\tagendprob

\begin{prop}[Kompaktlık]
  \ollabel{prop:proves-compact}
  \begin{enumerate}
  \item $\Gamma \Proves !A$ ise $\Gamma_0 \Proves !A$ olacak biçimde sonlu
    bir $\Gamma_0 \subseteq \Gamma$ alt kümesi vardır.
  \item $\Gamma$'nın her sonlu alt kümesi tutarlıysa $\Gamma$ tutarlıdır.
  \end{enumerate}
\end{prop}

\begin{proof}
  \begin{enumerate}
    \item $\Gamma \Proves !A$ ise $\Gamma$'dan $!A$ için
      !!a{derivation}~$\delta$ vardır. $\Gamma_0$, $\delta$'nın
      !!{undischarged} varsayımlarının kümesi olsun. Her !!{derivation} sonlu
      olduğundan $\Gamma_0$ yalnızca sonlu sayıda !!{sentence}s içerebilir.
      Dolayısıyla $\delta$, sonlu bir $\Gamma_0 \subseteq \Gamma$'dan
      $!A$ için !!a{derivation} olur.
    \item Bu, $!A \ident \lfalse$ özel durumu için (1)'in karşıt-tersidir.
  \end{enumerate}
\end{proof}

